\expandafter\ifx\csname CAPETINCLUDEMODE\endcsname\relax
\documentclass[11pt,a4paper]{article}
% preamble.tex — packages et commandes communes a toutes les fiches CAPET SII IE
\input{glyphtounicode}
\pdfgentounicode=1
\usepackage{cmap}
\usepackage[T1]{fontenc}
\usepackage[utf8]{inputenc}
\usepackage[french]{babel}
\usepackage[a4paper, top=1.6cm, bottom=1.6cm, left=1.4cm, right=1.4cm]{geometry}
\setlength{\headheight}{14.5pt}
\usepackage{amsmath, amssymb, amsfonts}
\usepackage{siunitx}
\sisetup{locale=FR, output-decimal-marker={,}, inter-unit-product=\ensuremath{{}\cdot{}}}
\usepackage{array, booktabs, multirow, tabularx}
\usepackage{xcolor}
\usepackage[most, breakable]{tcolorbox}
\usepackage{multicol}
\usepackage{enumitem}
\usepackage{mathtools}
\usepackage{lmodern}
\usepackage{microtype}
\usepackage{fancyhdr}
\usepackage{tikz}
\usetikzlibrary{positioning,arrows.meta}
\usepackage{lastpage}

% ── Couleurs par domaine ──────────────────────────────────────────────
\definecolor{colElec}{RGB}{26, 82, 148}       % Électrotechnique — bleu foncé
\definecolor{colElecLight}{RGB}{210, 225, 245} % fond bleu clair
\definecolor{colAuto}{RGB}{34, 120, 66}        % Automatique — vert
\definecolor{colMeca}{RGB}{185, 90, 20}        % Mécanique — orange
\definecolor{colTherm}{RGB}{160, 40, 40}       % Thermique — rouge
\definecolor{colQT}{RGB}{232, 244, 233}        % fond question type (vert très léger)
\definecolor{colJury}{RGB}{255, 243, 205}      % fond jury (jaune ambre)
\definecolor{colJuryFrame}{RGB}{200, 130, 0}   % bordure jury
\definecolor{colGrayLight}{RGB}{245, 245, 247} % fond gris neutre
\definecolor{colSeparator}{RGB}{180, 180, 180} % gris séparateur

% ── Style tcolorbox de section principale ────────────────────────────
\tcbset{
  ficheSection/.style={
    breakable,
    enhanced,
    colback=#1!7!white,
    colframe=#1,
    colbacktitle=#1,
    coltitle=white,
    fonttitle=\bfseries\normalsize,
    left=4pt, right=4pt, top=3pt, bottom=4pt,
    before skip=6pt, after skip=6pt,
  }
}

% ── Style encadré question type ──────────────────────────────────────
\tcbset{
  questiontype/.style={
    breakable,
    enhanced,
    colback=colQT,
    colframe=colElec!55!black,
    colbacktitle=colElec!55!black,
    coltitle=white,
    fonttitle=\bfseries\small,
    left=4pt, right=4pt, top=3pt, bottom=3pt,
    before skip=5pt, after skip=5pt,
  }
}

% ── Style encadré jury ───────────────────────────────────────────────
\tcbset{
  jury/.style={
    breakable,
    enhanced,
    colback=colJury,
    colframe=colJuryFrame,
    fonttitle=\bfseries\normalsize,
    left=4pt, right=4pt, top=3pt, bottom=3pt,
  }
}

% ── Mise en page ─────────────────────────────────────────────────────
\pagestyle{fancy}
\fancyhf{}
\renewcommand{\headrulewidth}{0.4pt}
\renewcommand{\footrulewidth}{0.3pt}

\setlist[itemize]{leftmargin=1.2em, itemsep=1pt, topsep=2pt, parsep=0pt}
\setlist[enumerate]{leftmargin=1.2em, itemsep=3pt, topsep=2pt, parsep=0pt}

\setlength{\parskip}{2pt}
\setlength{\columnsep}{10pt}

\def\CAPETBEGINFILE{\begin{document}}
\def\CAPETENDFILE{\end{document}}
\else
\def\CAPETBEGINFILE{}
\def\CAPETENDFILE{}
\fi

\fancyhead[L]{\textbf{\textcolor{colElec}{T14}} - Variation de vitesse et de couple des machines}
\fancyhead[R]{\small\textit{CAPET SII IE - Session 2026}}
\fancyfoot[C]{\small\thepage\,/\,\pageref{LastPage}}

\CAPETBEGINFILE

% -----------------------------------------------------------------------
%  BLOC EN-TeTE
% -----------------------------------------------------------------------

\begin{tcolorbox}[
  enhanced, breakable,
  colback=colElec,
  colframe=colElec,
  left=6pt, right=6pt, top=5pt, bottom=5pt,
  before skip=0pt, after skip=4pt
]
  {\color{white}\bfseries\large T14 \quad Variation de vitesse et de couple des machines}\\[2pt]
  {\color{colElec!30!white}\small
    \textbf{Domaine :} electrotechnique - Variateurs \quad
    \textbf{Presence :} \textbf{8/13} sujets (2013-2026) - theme \emph{tres frequent}
  }\\[1pt]
  {\color{colElec!20!white}\footnotesize
    Sous-themes : commande scalaire V/f - THD et commande MLI - variateur de frequence ATV630 - boucle de vitesse PI - freinage regeneratif 4 quadrants - ceblage variateur-automate
  }
\end{tcolorbox}

% -----------------------------------------------------------------------
%  1. RAPPELS DE COURS
% -----------------------------------------------------------------------

\begin{tcolorbox}[
  ficheSection=colElec,
  title={1. Rappels de cours},
]

% -- 1.1 Variateur de frequence ----------------------------------------
\textbf{1.1 Structure d'un variateur de frequence}

\smallskip
\begin{center}
\textbf{Reseau AC} $\xrightarrow{\text{Pont redresseur}}$ \textbf{Bus DC} $\xrightarrow{\text{Cellule de filtrage}}$ \textbf{Bus DC filtre} $\xrightarrow{\text{Onduleur MLI}}$ \textbf{MAS}
\end{center}
Avantages : contrele de la vitesse et du couple, demarrage progressif (rampe d'acceleration), freinage par injection de courant continu ou recuperatif, economie d'energie sur charges - couple variable ($P \propto N^3$).

\medskip\hrule\vspace{4pt}

% -- 1.2 Commande scalaire V/f ----------------------------------------
\textbf{1.2 Commande scalaire V/f}

\smallskip
\[
  \frac{V_s}{f_s} = \text{constante} \quad \Rightarrow \quad \text{flux magnetique constant}
\]
Principe : faire varier simultanement la frequence statorique $f_s$ et la tension $V_s$ pour maintenir le flux et donc le couple disponible.
Boucle de vitesse avec correcteur PI :
\[
  H_{BF}(p) = \frac{1}{1 + \dfrac{2m}{\omega_0}\,p + \dfrac{1}{\omega_0^2}\,p^2}
  \qquad;\qquad
  m = \text{amortissement}, \quad \omega_0 = \text{pulsation naturelle}
\]
Donnees 2024 : $K_p = 0{,}35$, $K_i = 2{,}5$, $m = 1$ (amortissement critique), erreur statique nulle, depassement nul.

\medskip\hrule\vspace{4pt}

% -- 1.3 THD et choix de commande ------------------------------------
\textbf{1.3 Taux de distorsion harmonique (THD) - commande pleine onde vs MLI}

\smallskip
\[
  THD = \frac{\sqrt{I_2^2 + I_3^2 + \cdots + I_n^2}}{I_1}
  \qquad;\qquad
  \text{Norme IEC/EN 61000-2-2 : } THD < 25\,\%
\]

\begin{center}
\renewcommand{\arraystretch}{1.25}
\begin{tabular}{lll}
\toprule
\textbf{Commande} & \textbf{THD typique} & \textbf{Conformite norme} \\
\midrule
Pleine onde & $\approx 59\,\%$ & Non (> 25 \%) \\
MLI ($f_p = 6\,\mathrm{kHz}$) & $\approx 25\,\%$ & Limite acceptable \\
\bottomrule
\end{tabular}
\end{center}
La MLI repousse les harmoniques vers les hautes frequences (autour de $f_p$ et ses multiples) - THD du fondamental ameliore.

\medskip\hrule\vspace{4pt}

% -- 1.4 Dimensionnement variateur ------------------------------------
\textbf{1.4 Dimensionnement et ceblage du variateur ATV630}

\smallskip
Courant nominal moteur : $I_n = \dfrac{P_u}{\eta\,\sqrt{3}\,U_{comp}\,\cos\varphi}$

\begin{center}
\renewcommand{\arraystretch}{1.25}
\begin{tabular}{lll}
\toprule
\textbf{Critere de choix} & \textbf{Parametre} & \textbf{Exemple ATV630C11N4} \\
\midrule
Puissance moteur & $P_u$ & 90\,kW / 100\,kW \\
Surcharge admissible & ND (110\,\%) ou HD (150\,\%) & HD si charge lourde \\
Tension reseau & 400\,V triphase & - \\
Pouvoir de coupure disjoncteur amont & $I_{cu}$ & 50\,kA (NSX250N MA220) \\
\bottomrule
\end{tabular}
\end{center}
\textbf{Ceblage :} sens de rotation (LI1 / LI2), reference de vitesse (entree 0-10\,V ou 4-20\,mA), retour de defaut - sortie DO. Rapport jury 2021 : maetrise du ceblage automate-variateur obligatoire.

\end{tcolorbox}

% -----------------------------------------------------------------------
%  2. FORMULAIRE SYNTHeTIQUE
% -----------------------------------------------------------------------

\begin{tcolorbox}[
  ficheSection=colElec,
  title={2. Formulaire synthetique},
]
\small
\begin{multicols}{2}
\renewcommand{\arraystretch}{1.45}
\begin{tabular}{@{}p{4.5cm}ll@{}}
\toprule
\textbf{Formule} & \textbf{Usage} & \textbf{U.} \\
\midrule
\multicolumn{3}{@{}l}{\footnotesize\textit{Commande V/f}} \\
$V_s/f_s = \mathrm{cste}$ & Loi scalaire & - \\
$N_s = 60f/p$ & Synchro & tr/min \\
\midrule
\multicolumn{3}{@{}l}{\footnotesize\textit{THD}} \\
$THD = \sqrt{\sum_{k\geq2} I_k^2}/I_1$ & Distorsion & - \\
$THD < 25\,\%$ & Norme & - \\
\midrule
\multicolumn{3}{@{}l}{\footnotesize\textit{Courant moteur}} \\
$I_n = P_u/(\eta\sqrt{3}U\cos\varphi)$ & Nominal & A \\
\bottomrule
\end{tabular}

\columnbreak

\renewcommand{\arraystretch}{1.45}
\begin{tabular}{@{}p{4.5cm}ll@{}}
\toprule
\textbf{Formule} & \textbf{Usage} & \textbf{U.} \\
\midrule
\multicolumn{3}{@{}l}{\footnotesize\textit{Boucle vitesse PI}} \\
$H(p) = K_p(1 + 1/(\tau_i p))$ & Correcteur & - \\
$m$ = amortissement & BF 2\textsuperscript{e} ordre & - \\
\midrule
\multicolumn{3}{@{}l}{\footnotesize\textit{Quadrants}} \\
I : $C>0,\Omega>0$ & Moteur direct & - \\
II : $C<0,\Omega>0$ & Freinage recup. & - \\
III,IV & Rotation inverse & - \\
\midrule
\multicolumn{3}{@{}l}{\footnotesize\textit{Donnees 2024}} \\
$f_p = 6\,\mathrm{kHz}$ & MLI & Hz \\
$THD_{pleine} = 59\,\%$ & Pleine onde & - \\
\bottomrule
\end{tabular}
\end{multicols}
\normalsize
\end{tcolorbox}

% -----------------------------------------------------------------------
%  3. QUESTIONS TYPES
% -----------------------------------------------------------------------

\begin{tcolorbox}[
  ficheSection=colElec,
  title={3. Questions types - extraites des sujets reels},
]

% -- QT1 -------------------------------------------------------------

% Mini-procedure de resolution type copie concours
\begin{tcolorbox}[questiontype,title={Procedure type copie concours (5 etapes)}]
\begin{enumerate}[leftmargin=1.5em]
  \item \textbf{Qualifier le modele} : hypotheses, domaine de validite, grandeurs d'entree/sortie.
  \item \textbf{Ecrire les lois} : equations de base avec symboles et unites SI.
  \item \textbf{Isoler l'inconnue} : expression litterale, puis application numerique tracee.
  \item \textbf{Controler l'ordre de grandeur} : coherence physique et unite finale.
  \item \textbf{Conclure concours} : phrase courte orientee decision technique.
\end{enumerate}
\end{tcolorbox}

% Visuel de methode : chemin de resolution CAPET
\begin{center}
\begin{tikzpicture}[font=\small, node distance=7mm and 6mm, >=Stealth]
  \node[draw, rounded corners, fill=colElec!8, inner sep=4pt] (n1) {Notions};
  \node[draw, rounded corners, fill=colElec!8, right=of n1, inner sep=4pt] (n2) {Methode};
  \node[draw, rounded corners, fill=colElec!8, right=of n2, inner sep=4pt] (n3) {Application annale};
  \node[draw, rounded corners, fill=colElec!8, right=of n3, inner sep=4pt] (n4) {Verification jury};
  \draw[->, thick, colElec] (n1) -- (n2);
  \draw[->, thick, colElec] (n2) -- (n3);
  \draw[->, thick, colElec] (n3) -- (n4);
\end{tikzpicture}
\end{center}
\begin{tcolorbox}[questiontype,
  title={QT1 - THD et conformite norme - commande MLI vs pleine onde \hfill\normalfont\footnotesize E2 disc. 2024 (B10-B11)}
]
\textbf{enonce.}
Le DT3 donne le spectre du courant reseau en pleine onde : $I_1 = 170\,\mathrm{A}$, $I_3 = 95\,\mathrm{A}$, $I_5 = 27\,\mathrm{A}$, $I_7 = 27\,\mathrm{A}$, $I_9 = 20\,\mathrm{A}$. En MLI ($f_p = 6\,\mathrm{kHz}$) : $I_1 = 170\,\mathrm{A}$, $I_3 = 37\,\mathrm{A}$, $I_5 = 22\,\mathrm{A}$. Calculer le THD dans chaque cas et conclure.

\tcbline

\textbf{Corrige.}

\textbf{Pleine onde :}
\[
  THD = \frac{\sqrt{95^2 + 27^2 + 27^2 + 20^2}}{170} = \frac{\sqrt{9025 + 729 + 729 + 400}}{170} = \frac{\sqrt{10883}}{170} = \frac{104{,}3}{170} \approx \mathbf{61{,}4\,\%}
\]
Non conforme (> 25\,\%).

\textbf{MLI :}
\[
  THD = \frac{\sqrt{37^2 + 22^2}}{170} = \frac{\sqrt{1369 + 484}}{170} = \frac{\sqrt{1853}}{170} = \frac{43{,}0}{170} \approx \mathbf{25{,}3\,\%}
\]
Conforme ou limite ($\leq 25\,\%$). La MLI est necessaire pour respecter la norme IEC/EN 61000-2-2.
\end{tcolorbox}

\vspace{4pt}

% -- QT2 -------------------------------------------------------------

\begin{tcolorbox}[questiontype,
  title={QT2 - Choix et dimensionnement du variateur ATV630 \hfill\normalfont\footnotesize Rapport jury 2021 (Q15-Q18)}
]
\textbf{enonce.}
Choisir le modele ATV630 pour un moteur de $P_u = 90\,\mathrm{kW}$, $U = 400\,\mathrm{V}$, $\cos\varphi = 0{,}86$, $\eta = 0{,}942$, surcharge 150\,\% (mode HD). Indiquer le disjoncteur amont ($I_{cu,reseau} = 50\,\mathrm{kA}$).

\tcbline

\textbf{Corrige.}

\textbf{Courant nominal moteur :}
\[
  I_n = \frac{P_u}{\eta\,\sqrt{3}\,U\,\cos\varphi} = \frac{90\,000}{0{,}942 \times 1{,}732 \times 400 \times 0{,}86} \approx \frac{90\,000}{562} \approx 160\,\mathrm{A}
\]

\textbf{Choix variateur :} ATV630\textbf{C11N4} ($I_n = 170\,\mathrm{A}$, 110\,kW, HD 150\,\%, 400\,V). La puissance 110\,kW > 90\,kW et le courant 170\,A > 160\,A.

\textbf{Disjoncteur amont :} NSX250N MA220 - declencheur \textbf{MA} (magnetique seul, seuil $I_{rm} = 2860\,\mathrm{A}$, le variateur assure la protection thermique), $I_{cu} = 50\,\mathrm{kA}$ (adapte au $I_{cc}$ reseau).
\end{tcolorbox}

\vspace{4pt}

% -- QT3 -------------------------------------------------------------

\begin{tcolorbox}[questiontype,
  title={QT3 - Boucle de vitesse scalaire - correcteur PI \hfill\normalfont\footnotesize E2 disc. 2024 (C5-C11)}
]
\textbf{enonce.}
La boucle de vitesse de la MAS est un second ordre avec $K_{meca} = 250$, $\tau_{meca} = 8{,}5\,\mathrm{s}$. On souhaite un amortissement $m = 1$ (aperiodique) et un temps de reponse $T_r = 1\,\mathrm{s}$. Exprimer $K_p$ et $K_i$.

\tcbline

\textbf{Corrige.}

La fonction de transfert en boucle fermee avec correcteur PI est du second ordre :
\[
  H_{BF}(p) = \frac{1}{1 + \frac{2m}{\omega_0}\,p + \frac{1}{\omega_0^2}\,p^2}
\]
Avec $m = 1$ et $\omega_0 = 2\pi/T_r \times 2m$ (critere amortissement) :
\[
  \omega_0 = \frac{1}{\tau_i K_p K_{meca}} \quad;\quad 2m = \frac{K_p K_{meca}\tau_{meca} + \tau_i}{\tau_i \omega_0}
\]
Application numerique (resultat donne dans le sujet) : $K_p \approx 0{,}35$, $K_i \approx 2{,}5$.

\textbf{Verification :} depassement nul ($m = 1$), erreur statique nulle (integrateur dans le PI), marge de phase $= 180^\circ$.
\end{tcolorbox}

\end{tcolorbox}

% -----------------------------------------------------------------------
%  4. ATTENTES DU JURY
% -----------------------------------------------------------------------

\begin{tcolorbox}[
  jury,
  title={4. Attentes du jury - Erreurs recurrentes et points de vigilance},
]

\begin{enumerate}[label=\textbf{\arabic*.}, leftmargin=*, itemsep=5pt]

  \item \textbf{Justifier le choix du variateur - pas seulement la reference.}
  Citer les criteres : puissance ($P_u$), type de service (ND ou HD), tension reseau, pouvoir de coupure. Rapport jury 2021 : "sans justification, sans indiquer clairement les criteres de choix".

  \item \textbf{Disjoncteur MA devant un variateur.}
  Le variateur assure lui-meme la protection thermique du moteur. Utiliser un declencheur magnetique seul (MA) en amont. Un disjoncteur magneto-thermique cree une double protection thermique - decoordination.

  \item \textbf{THD : mettre au carre sous la racine.}
  Oublier de mettre les harmoniques au carre, ou diviser par l'amplitude totale au lieu de $I_1$, sont les erreurs les plus frequentes.

  \item \textbf{Commande scalaire vs vectorielle.}
  La commande scalaire V/f suffit pour les charges - couple variable (pompe, ventilateur). La commande vectorielle FOC n'est pas exigee au CAPET mais peut etre citee pour les applications - dynamique elevee.

  \item \textbf{Ceblage automate-variateur.}
  Rapport jury 2021 : "un candidat au CAPET doit etre capable de cebler une entree ou une sortie TOR-analogique d'un automate industriel programmable, ainsi que la partie commande d'un variateur de frequence".

  \item \textbf{Freinage regeneratif et 4 quadrants.}
  Identifier les quadrants (couple/vitesse). Un onduleur 4 quadrants est necessaire pour le freinage recuperatif (quadrant II ou IV). Verifier que la structure materielle le permet.

\end{enumerate}

\medskip
\begin{tcolorbox}[
  colback=colJuryFrame!12!colJury,
  colframe=colJuryFrame,
  left=4pt, right=4pt, top=2pt, bottom=2pt,
  fontupper=\small\itshape
]
Rapport jury 2021 : "Le jury constate que beaucoup de candidats ne sont pas - l'aise dans la realisation de schemas electriques. Un candidat au CAPET doit etre capable de cebler une entree ou une sortie TOR-analogique d'un automate industriel programmable, ainsi que la partie commande d'un variateur de frequence."\\
Rapport jury 2024 : "Le contrele de la vitesse peut etre realise via une commande scalaire de la machine asynchrone et une boucle d'asservissement. L'objectif est d'etudier les parametres - respecter : erreur statique, depassement et marge de phase."
\end{tcolorbox}

\end{tcolorbox}

% -----------------------------------------------------------------------
%  PIED - PReSENCE DANS LES SUJETS
% -----------------------------------------------------------------------

\begin{tcolorbox}[
  enhanced, breakable,
  colback=colGrayLight,
  colframe=colSeparator,
  left=4pt, right=4pt, top=3pt, bottom=3pt,
  before skip=4pt, after skip=0pt,
  fonttitle=\small\bfseries,
  title={Presence dans les sujets analyses (2013-2026)},
]
\begin{center}
\footnotesize\renewcommand{\arraystretch}{1.15}
\begin{tabular}{lccccccccccccc}
\toprule
\textbf{Annee} & 13 & 14 & 15 & 16 & 17 & 18 & 19 & 20 & 21 & 22 & 23 & 24 & 25 \\
\midrule
E1 (commune)   & - & - & - & - & - & - & - & $\checkmark$ & - & - & - & $\checkmark$ & - \\
E2 discipl. IE & - & - & - & $\checkmark$ & - & $\checkmark$ & - & - & $\checkmark$ & - & - & $\checkmark$ & - \\
E2 appl. comm. & - & - & - & - & - & - & - & - & - & - & - & - & - \\
\bottomrule
\end{tabular}
\end{center}
\smallskip
\footnotesize\textbf{Sous-themes frequents :}
commande V/f et THD (4/13) - choix/ceblage variateur ATV (3/13) - boucle de vitesse PI (2/13) - freinage 4 quadrants (3/13) - dimensionnement disjoncteur MA (3/13)

\medskip
\textbf{Ordres de grandeur (reperes rapides).}
\begin{itemize}[leftmargin=1.3em]
  \item Reseau BT FR : tension phase-neutre typique \SI{230}{\volt}, tension phase-phase typique \SI{400}{\volt}, frequence \SI{50}{\hertz}.
  \item Courants d'etude frequents en sujet : de quelques \si{\ampere} a quelques dizaines de \si{\ampere}.
  \item En conversion d'energie, un rendement attendu reste generalement entre 0{,}8 et 0{,}98 selon la technologie.
\end{itemize}

\medskip
\textbf{Synthese express (a refaire sans support).}
\begin{enumerate}[leftmargin=1.4em]
  \item Poser les hypotheses du modele et les conventions de signe.
  \item Ecrire 2 a 4 relations cles avec unites.
  \item Realiser un calcul numerique complet et verifier l'ordre de grandeur.
  \item Justifier le resultat par un argument technique de niveau CAPET.
\end{enumerate}

\end{tcolorbox}

\CAPETENDFILE









